\chapter{The stationary state}\label{ch:stationary}

A drop adds one grain and the avalanche that follows removes some; over a long run the
board settles into a state in which macroscopic quantities stop drifting while the
configuration keeps changing. This chapter describes that state: which configurations
it consists of, what the mean height is, and what is known and expected about the
distributions of the four avalanche observables of \eqref{eq:observables}.

\section{Recurrent configurations and the uniform measure}\label{sec:recurrent}

Not every stable configuration is reachable once the board has been driven for a while.
The states the dynamics returns to form the set of \emph{recurrent} configurations
$\mathcal{R}$; the rest are transient in the sense that after enough drops they are
never visited again. On $\mathcal{R}$ the two ways of sampling a system agree.
Sampling \emph{over time} means recording the configurations visited by one long run;
sampling an \emph{ensemble} means drawing them independently from the stationary
distribution. These differ while the initial condition still influences the run, which
is what the burn-in of §\ref{sec:burnin} removes, and they differ if the chain is not
ergodic. Drop-and-relax is irreducible and aperiodic on $\mathcal{R}$ --- every
recurrent configuration is reachable from every other, and the chain does not cycle ---
so the stationary distribution there is unique, and the time average along a single run
converges to the ensemble average.

That measure is the uniform one. Every recurrent configuration carries the same weight
$1/\absR$, and the number of them is
\begin{equation}\label{eq:recurrentcount}
  \absR = \det(-\D),
\end{equation}
by the matrix--tree theorem, since $\det(-\D)$ counts the spanning trees of the wired
graph obtained by attaching a sink to every off-board direction
\cite{dhar1999,jarai2018,klivans2019}. The uniformity is what makes the mean height an
exactly computable number rather than a fitted one, and it is the reason a measurement
of $\meanh$ is a sharp check of the whole implementation.

The set $\mathcal{R}$ also carries an algebraic structure: the addition operators
generate a finite abelian group acting on it, the sand pile group, whose order is
$\absR$ \cite{dhar1995}. Beyond the count \eqref{eq:recurrentcount} nothing in this
report depends on that structure, but the group order is computed separately in
Chapter~\ref{ch:hexagonal}.

\section{The mean height}\label{sec:meanheight}

The mean height is the average of the height function over the stationary measure,
\begin{equation}\label{eq:meandef}
  \meanh = \frac{1}{L^2}\sum_{\bm{x}} p(\bm{x}) ,
\end{equation}
where the average is taken both over sites and over the ensemble. A first guess is that
the four stable heights $0,1,2,3$ are equally likely, giving $1.5$. That guess ignores a
constraint that the recurrent configurations carry: two neighbouring sites cannot both
be empty. The constraint comes from the burning algorithm, in which a site of height
zero can be burnt only after all of its neighbours, so that neither of two adjacent
empty sites could ever be burnt and the configuration would not be recurrent. Counting
the configurations subject to that restriction gives the exact value
\begin{equation}\label{eq:meanh}
  \meanh = \frac{17}{8} = 2.125 ,
\end{equation}
which is already well above the uniform guess.

The individual height probabilities are known exactly as well. Priezzhev computed them
for the square lattice \cite{priezzhev1994}, and the closed forms were later reduced by
Jeng, Piroux and Ruelle \cite{jeng2006} to
\begin{align}
  p(0) &= \tfrac{2}{\pi^2} - \tfrac{4}{\pi^3} = 0.0736, &
  p(1) &= \tfrac{1}{4} - \tfrac{1}{2\pi} - \tfrac{3}{\pi^2} + \tfrac{12}{\pi^3} = 0.1739,
  \notag \\[2pt]
  p(2) &= \tfrac{3}{8} + \tfrac{1}{\pi} - \tfrac{12}{\pi^3} = 0.3063, &
  p(3) &= \tfrac{3}{8} - \tfrac{1}{2\pi} + \tfrac{1}{\pi^2} + \tfrac{4}{\pi^3} = 0.4462 ,
  \label{eq:heightprobs}
\end{align}
and $\sum_i i\,p(i)$ reproduces $2.125$; the single integral those formulas contain was
later evaluated and proved to take the value that makes this true
\cite{poghosyan2011,levine2014}. The four probabilities and the mean are the sharpest
available test of the simulator: they involve no fitting, no cut-off and no assumption
about a tail, so a disagreement is a defect rather than a shortage of samples. They also
give the exact relation between the threshold and the mean height that
Chapter~\ref{ch:hexagonal} needs, since they show that $\meanh$ is a property of the
lattice and not a constant.

\section{The avalanche distributions}\label{sec:distributions}

For each avalanche the record holds the four integers of \eqref{eq:observables}, and the
question is what distribution each of them follows. In the power-law regime the
complementary cumulative distribution of a size $X$ behaves as
\begin{equation}\label{eq:ccdf}
  P(X \ge x) \sim x^{-(\tau-1)} ,
\end{equation}
so the slope of the straight part on log--log axes gives $\tau - 1$. The exponent is
better estimated by maximum likelihood than by fitting a histogram, since a histogram
depends on the binning and the likelihood does not:
\begin{equation}\label{eq:mle}
  \hat{\tau} = 1 + \frac{n}{\sum_i \ln(x_i / x_{\min})},
\end{equation}
where the sum runs over the $n$ events with $x_i \ge x_{\min}$. The two conventions ---
probability density $x^{-\tau}$ and complementary cumulative $x^{-(\tau-1)}$ --- differ
by one, and a table of exponents is meaningless unless the convention is stated.

For the toppling number the model has a conjectured exact exponent. The
two-dimensional undirected sand pile is equivalent to the $q \to 0$ limit of the Potts
model \cite{majumdar1992}, and the exponents read off from that equivalence are
$\tau_s = 6/5$, $\tau_d = 5/4$ and $\tau_t = 7/5$ for topplings, distinct sites and
duration, in the probability-density convention \cite{dhar1999,ivashkevich1998}. In the
cumulative convention of \eqref{eq:ccdf} these correspond to slopes $0.2$ and $0.25$,
which is why the distributions are so flat that they can be mistaken for a cut-off over
the range a small board can reach.

Two properties of the measurement follow from the finiteness of the board and belong in
the analysis from the start. An avalanche cannot exceed the board, so the distribution is cut off at a scale set by $L$, and an estimate of the
exponent taken with a lower cut-off near that scale describes the steep edge of the
cut-off rather than the power law, returning a value that is too high. The remedy is not
a better estimator but a range of sizes: the cumulative distributions of several boards
are plotted together, and the exponent is read from a straight section that is common to
all of them and lies before the cut. Second, the length $R$ is a geometric quantity
bounded by the board and is not expected to be a pure power law at all; it belongs in a
finite-size collapse, $s/N^{D}$ for a cut-off exponent $D$, rather than in a fit of its
own. If the collapsed curves fail to overlap, $D$ is the first thing to adjust.

\begin{pending}
  \pendingtag The exponents $\tau_s$ and $\tau_d$ have not been measured. Doing so
  requires a run whose usable tail holds on the order of $10^3$ events or more at each
  of at least three board sizes, so that the collapse described above can be checked
  before any exponent is quoted.
\end{pending}

\section{Correlations between the four observables}\label{sec:correlations}

The four observables are not independent, and the expected relations follow from their
definitions in \eqref{eq:observables}.

\emph{Topples against area.} The area counts sites toppled at least once and the toppling
number counts all topplings, so $s \ge s_d$, with equality exactly when no site topples
twice. In a small avalanche that is the typical case and the two are proportional; in a
large avalanche sites near the drop point topple repeatedly, so $s$ grows faster than
$s_d$ and one expects $s \sim s_d^{\gamma}$ with $\gamma > 1$. The relation is visible
directly as a rising mean of the ratio $s/s_d$ with $s_d$.

\emph{Area against length.} The area of an avalanche is related to its radius by a
fractal dimension, $s_d \sim R^{D_f}$. If avalanches filled their bounding region then
$D_f$ would be $2$; a value below $2$ means the cluster is ramified, with a perimeter
comparable to the area.

\emph{Loss against length.} This relation is a threshold, not a power law, and it is the
sharpest of the four. Loss requires a boundary site to topple, so $l = 0$ for every
avalanche that does not reach the boundary; once $R$ exceeds the distance from the drop
site to the nearest boundary, loss begins and grows quickly with $R$. A linear
correlation coefficient is the wrong summary here, because it averages the exact zeros
below the threshold with the growing values above it and reports a weak relation where
the relation is in fact deterministic in shape. The appropriate summary is the
conditional mean $E[l \mid R]$ after binning by $R$, which reproduces the threshold as a
corner at the boundary distance.

\emph{Loss against topples.} Weak, for the same reason: most topplings are interior and
conserve sand, so a large avalanche that stays away from the boundary loses nothing.

\begin{pending}
  \pendingtag The correlation coefficients and the binned conditional mean $E[l \mid R]$
  have not been computed. The recorded columns are sufficient for both; the work is
  post-processing of the avalanche record rather than further simulation.
\end{pending}
